{
	\centering
	\pgfplotstableread[col sep=comma] {../misc/short_force_perturbations.csv}\datatable
	%\resizebox{8.69cm}{!}{	
		\tikzsetnextfilename{fig_short_perturbations_implemented}		
		\begin{tikzpicture}
			\centering
			%
			\begin{axis}[%
				name=ax1,
				%width=.48\linewidth,
				width=4cm,
				height=3cm,
				scale only axis,
				xlabel={\footnotesize{Time (\si{\milli\second})}},
				ylabel={\footnotesize{Force (\si{\newton})}},
				axis x line = bottom,
				axis y line = left,
				ylabel near ticks,
				ylabel style={yshift=-.4cm},
				xlabel style={yshift=1cm},
				xticklabel style = {font=\footnotesize},
				yticklabel style = {font=\footnotesize},
				]
				%
				\addplot [
					color=Orient,
					smooth,
					thick,
					]
					table [
					col sep=comma, 
					x expr=\coordindex, 
					y=force_pert_pos
					] 
					{\datatable};
				\addplot [
					color=Shiraz,
					smooth,
					mark=none,
					thick,
					]
					table [
					col sep=comma, 
					x expr=\coordindex, 
					y=force_pert_neg] 
					{\datatable};
			\end{axis}
			%
%			\begin{axis}[%
%				width=.48\linewidth,
%				height=.2\linewidth,
%				scale only axis,
%				xlabel={time (\si{\milli\second})},
%				%ylabel={Force (\SI{}{\newton})},
%				axis x line = bottom,
%				axis y line = left, 			
%				yticklabel style={
%					%/pgf/number format/precision=3,
%					/pgf/number format/fixed},
%				at={($(ax1.south east)+(6,0)$)},
%				%scaled y ticks=false,
%				]
%				%
%				\addplot [
%				color=Orient,
%				smooth,
%				mark=none,
%				thick,
%				]
%				table [col sep=comma, x expr=\coordindex, y=force_pert_neg] {\datatable};
%			\end{axis}
			
		\end{tikzpicture}
%	}
	%\caption[Exemples de perturbations]{Exemples d'une perturbation positive et d'une négative, qui ont pu être extraites en utilisant la méthode décrite en \cref{chap3_sec_method_sub_virt_traj_6}.}
	%\label{fig_short_perturbations_implemented} % Label 
}